% ============================================================
% fig_inversa_tres_roc.tex
% Las tres antitransformadas de X(s) = (s+3)/((s+1)(s+2)) según la ROC
% (ejemplo de §7.3.2). Cada fila: [plano s con ROC] L^{-1} [x(t)].
%   (a) ROC Re{s} > -1:        x(t) = (2e^{-t} - e^{-2t}) u(t)   causal
%   (b) ROC Re{s} < -2:        x(t) = (-2e^{-t} + e^{-2t}) u(-t) anticausal
%   (c) ROC -2 < Re{s} < -1:   x(t) = -2e^{-t}u(-t) - e^{-2t}u(t) bilateral
% Estilo consistente con fig_polo_real_multiple y fig_roc_tres_formas.
% Las curvas temporales están reescaladas (sketch cualitativo).
% ============================================================
%\PassOptionsToPackage{gray}{xcolor}

\documentclass[border=5mm]{standalone}
% ==============================================================================
% SETUP.TEX - CONFIGURACIÓN GLOBAL DEL PROYECTO
% ==============================================================================
% Este archivo centraliza todos los paquetes y estilos.
% Debe incluirse en el preámbulo del libro principal y de los archivos standalone.
% \PassOptionsToPackage{gray}{xcolor}
% --- 1. CODIFICACIÓN E IDIOMA ---
% fix-cm habilita las fuentes Computer Modern en tamaños arbitrarios (por debajo
% de 5 pt, entre otros). Debe cargarse antes que fontenc.
\usepackage{fix-cm}
\usepackage[utf8]{inputenc}
\usepackage[T1]{fontenc}
\usepackage[spanish, es-tabla]{babel}  
\usepackage{csquotes} % Recomendado para babel y citas

\usepackage{import}
% mode=image: Usa el PDF pre-compilado, nunca intenta recompilar.
% Debes compilar cada figura manualmente antes de compilar el libro.
\usepackage[mode=image, subpreambles=false]{standalone}

% --- 2. MATEMÁTICAS Y CIENCIAS ---
\usepackage{amsmath, amssymb, amsfonts, mathtools}
\usepackage{enumitem}

% LaTeX trae \DeclareMathSizes{5}{5}{5}{5}, así que a 5 pt (\tiny) los subíndices
% salen del mismo tamaño que la base: en las etiquetas de figuras el M de
% $\omega_M$ queda desproporcionado. Se restablece la proporción habitual.
\DeclareMathSizes{5}{5}{3.7}{3}

% --- 3. DISEÑO Y GEOMETRÍA --- 
% Unificamos los márgenes para todo el documento
\usepackage[left=2.5cm,right=2.5cm,top=3cm,bottom=3cm]{geometry}
\makeatletter
\ifstandalone % Evita que geometry dañe el recorte de las figuras bajándolas 3cm
  \@ifclasswith{standalone}{crop=false}{
    % Es un capítulo/sección: Dejamos que geometry actúe.
  }{
    % Es una figura: Neutralizamos geometry para que no las recorte mal.
    \geometry{pass}
  }
\fi
\makeatother
\usepackage{parskip} % Espaciado entre párrafos sin sangría
\usepackage{float}   % Para usar [H] en figura

% Ajustes de flotantes para evitar espacios verticales exagerados
\renewcommand{\topfraction}{0.95}      % Permite que las figuras ocupen hasta el 95% superior
\renewcommand{\bottomfraction}{0.95}   % Permite que las figuras ocupen hasta el 95% inferior
\renewcommand{\textfraction}{0.05}     % Permite hojas con tan solo un 5% de texto
\renewcommand{\floatpagefraction}{0.8} % Requiere que una página de solo figuras esté 80% llena
\setcounter{topnumber}{4}
\setcounter{bottomnumber}{4}
\setcounter{totalnumber}{10}

% --- 4. GRÁFICOS Y TABLAS ---
\usepackage{graphicx}
\usepackage{booktabs} % Tablas profesionales
\usepackage{caption}  % Control de captions y \captionof
\usepackage{tikz}
\usetikzlibrary{arrows.meta, calc, angles, quotes, babel, backgrounds}
\usepackage{pgfplots}
\usepgfplotslibrary{groupplots}
\usepgfplotslibrary{external}
\usepgfplotslibrary{patchplots}
\pgfplotsset{compat=1.18}
\usepackage[american, straightvoltages]{circuitikz}

% --- 5. COLORES Y CAJAS ---

\usepackage[table]{xcolor}
\usepackage{tcolorbox}

% Definición de colores corporativos/estilo
\definecolor{utnblue}{RGB}{0, 51, 102}
\definecolor{codegreen}{rgb}{0,0.6,0}
\definecolor{codegray}{rgb}{0.5,0.5,0.5}
\definecolor{codepurple}{rgb}{0.58,0,0.82}
\definecolor{backcolour}{rgb}{0.96,0.96,0.96}

% --- 6. ENTORNOS PERSONALIZADOS (CAJAS) ---

% Caja "Resultado" (Azul Corporativo) — Definiciones, fórmulas clave, resultados importantes.
% Uso: \begin{resultado}[Fórmula de Euler] ... \end{resultado}
\newtcolorbox{resultado}[1][]{
    colback=utnblue!5!white,
    colframe=utnblue,
    title=\textbf{#1},
    fonttitle=\bfseries,
    sharp corners,
    boxrule=0.5mm
}

% Caja "Nota" (Gris) — Advertencias, aplicaciones, contexto secundario.
% Uso: \begin{nota}[Cuidado con la calculadora] ... \end{nota}
% Uso: \begin{nota} ... \end{nota}  → título por defecto "Nota"
\newtcolorbox{nota}[1][Nota]{
    colback=codegray!5!white,
    colframe=codegray,
    title=\textbf{#1},
    fonttitle=\bfseries,
    sharp corners,
    boxrule=0.5mm
}

% Caja inline para resaltar un valor y enlazarlo con su otra aparición en una tabla
% o en una cuenta: mismo color, mismo valor.
% Uso: \marcavalor{$4$}  o  \marcavalor[codegreen]{$4$}
\newtcbox{\marcavalor}[1][utnblue]{
    on line,
    boxsep=1pt, left=2pt, right=2pt, top=1pt, bottom=1pt,
    colback=#1!10!white,
    colframe=#1,
    boxrule=0.4pt,
    arc=2pt
}

% --- 7. CÓDIGO FUENTE (LISTINGS) ---
\usepackage{listings}

\lstdefinestyle{miestilo}{
    backgroundcolor=\color{backcolour},   
    commentstyle=\color{codegreen},
    keywordstyle=\color{magenta},
    numberstyle=\tiny\color{codegray},
    stringstyle=\color{codepurple},
    basicstyle=\ttfamily\footnotesize,
    breakatwhitespace=false,         
    breaklines=true,                 
    captionpos=b,                    
    keepspaces=true,                 
    numbers=left,                    
    numbersep=5pt,                  
    showspaces=false,                
    showstringspaces=false,
    showtabs=false,                  
    tabsize=2,
    frame=single,                    
    rulecolor=\color{codegray}
}

\lstset{style=miestilo}
% Soporte para tildes y ñ en bloques de código
\lstset{literate={ñ}{{\~n}}1 {á}{{\'a}}1 {é}{{\'e}}1 {í}{{\'i}}1 {ó}{{\'o}}1 {ú}{{\'u}}1}

% Operadores matemáticos personalizados

% Vector en sentido discreto (lista finita de coordenadas): negrita + flecha.
% Uso: \vect{v}, \vect{e}_k, \vect{0}.
\newcommand{\vect}[1]{\vec{\mathbf{#1}}}

% Fecha de versión y comando para capítulos standalone
\providecommand{\verdate}{\the\year.\ifnum\month<10 0\fi\the\month.\ifnum\day<10 0\fi\the\day}
\newcommand{\chapterversion}{%
    \vspace{-1.5cm}\begin{flushright}\small\textit{\color{codegray}Versión~\verdate}\end{flushright}\vspace{0.5cm}%
}

% --- 8. HIPERVÍNCULOS (DEBE SER EL ÚLTIMO PAQUETE) ---
\usepackage[colorlinks=true, linkcolor=utnblue, urlcolor=utnblue, citecolor=utnblue]{hyperref}

% --- 9. REFERENCIAS CRUZADAS ENTRE CAPÍTULOS STANDALONE ---
% Cuando se compila un capítulo standalone, xr-hyper lee los .aux de los
% otros capítulos (si existen) para resolver \ref{} a labels externos.
% En la compilación del libro completo este bloque no se activa.
\ifstandalone
  \usepackage{xr-hyper}
  \makeatletter
  \newcommand{\xrchap}[1]{%
    \edef\xrc@self{\detokenize{Capitulo#1}}%
    \edef\xrc@job{\jobname}%
    \ifx\xrc@self\xrc@job\else
      \IfFileExists{../#1capitulo/Capitulo#1.aux}%
        {\externaldocument{../#1capitulo/Capitulo#1}}{}%
    \fi
  }
  \makeatother
  \xrchap{01}\xrchap{02}\xrchap{03}\xrchap{04}
  \xrchap{05}\xrchap{06}\xrchap{07}\xrchap{08}
  \xrchap{09}\xrchap{10}\xrchap{11}\xrchap{12}
  \xrchap{13}
\fi

\begin{document}
\begin{tikzpicture}[>=Latex, scale=0.85, font=\small]

% =====================================================================
% Caso (a): ROC a la derecha de ambos polos — señal causal.
% =====================================================================
\begin{scope}[shift={(0, 0)}]
    % --- Plano s ---
    \fill[cyan!15] (-1, -1.7) rectangle (1.5, 1.7);
    \draw[->, thick] (-3, 0) -- (1.7, 0) node[right, font=\footnotesize] {$\sigma$};
    \draw[->, thick] (0, -1.7) -- (0, 1.9) node[above, font=\footnotesize] {$j\omega$};
    \draw[dashed, thick, gray] (-1, -1.7) -- (-1, 1.7);

    \node[red, font=\normalsize] at (-1, 0) {$\times$};
    \node[red, font=\normalsize] at (-2, 0) {$\times$};
    \node[below, font=\scriptsize] at (-1, -0.12) {$-1$};
    \node[below, font=\scriptsize] at (-2, -0.12) {$-2$};

    \node[utnblue!80!black, font=\scriptsize] at (0.9, 1.35) {ROC};

    % --- Conector L^{-1} ---
    \draw[->, gray, thick] (1.9, 0) -- (2.9, 0)
        node[midway, above, font=\scriptsize, gray] {$\mathcal{L}^{-1}$};

    % --- x(t): causal, arranca en 1 y decae ---
    \begin{scope}[shift={(3.2, 0)}]
        \draw[->, thick] (-0.5, 0) -- (3.4, 0) node[right, font=\footnotesize] {$t$};
        \draw[->, thick] (0, -0.4) -- (0, 1.7) node[right, font=\footnotesize] {$x(t)$};

        % 2e^{-t} - e^{-2t}: vale 1 en t = 0 (dibujada con factor 1.2)
        \draw[utnblue!90!black, very thick, smooth, samples=80, domain=0:3.2]
            plot (\x, {1.2*(2*exp(-\x) - exp(-2*\x))});
    \end{scope}

    \node[anchor=south, font=\small] at (2.0, 2.1) {(a)};
\end{scope}

% =====================================================================
% Caso (b): ROC a la izquierda de ambos polos — señal anticausal.
% =====================================================================
\begin{scope}[shift={(0, -5.4)}]
    % --- Plano s ---
    \fill[cyan!15] (-3, -1.7) rectangle (-2, 1.7);
    \draw[->, thick] (-3, 0) -- (1.7, 0) node[right, font=\footnotesize] {$\sigma$};
    \draw[->, thick] (0, -1.7) -- (0, 1.9) node[above, font=\footnotesize] {$j\omega$};
    \draw[dashed, thick, gray] (-2, -1.7) -- (-2, 1.7);

    \node[red, font=\normalsize] at (-1, 0) {$\times$};
    \node[red, font=\normalsize] at (-2, 0) {$\times$};
    \node[below, font=\scriptsize] at (-1, -0.12) {$-1$};
    \node[below, font=\scriptsize] at (-2, -0.12) {$-2$};

    \node[utnblue!80!black, font=\scriptsize] at (-2.5, 1.35) {ROC};

    % --- Conector L^{-1} ---
    \draw[->, gray, thick] (1.9, 0) -- (2.9, 0)
        node[midway, above, font=\scriptsize, gray] {$\mathcal{L}^{-1}$};

    % --- x(t): anticausal, crece hacia t -> -inf, dip negativo cerca de 0 ---
    \begin{scope}[shift={(6.4, 0)}]
        \draw[->, thick] (-3.3, 0) -- (0.5, 0) node[right, font=\footnotesize] {$t$};
        \draw[->, thick] (0, -0.5) -- (0, 1.7) node[right, font=\footnotesize] {$x(t)$};

        % (-2e^{-t} + e^{-2t}) sobre t < 0, factor 0.22; cruza cero en t = -ln 2
        \draw[utnblue!90!black, very thick, smooth, samples=80, domain=-1.3:0]
            plot (\x, {0.22*(-2*exp(-\x) + exp(-2*\x))});
        \draw[utnblue!90!black, very thick, dashed, smooth, samples=20, domain=-1.37:-1.3]
            plot (\x, {0.22*(-2*exp(-\x) + exp(-2*\x))});
    \end{scope}

    \node[anchor=south, font=\small] at (2.0, 2.1) {(b)};
\end{scope}

% =====================================================================
% Caso (c): ROC franja entre los polos — señal bilateral.
% =====================================================================
\begin{scope}[shift={(0, -10.8)}]
    % --- Plano s ---
    \fill[cyan!15] (-2, -1.7) rectangle (-1, 1.7);
    \draw[->, thick] (-3, 0) -- (1.7, 0) node[right, font=\footnotesize] {$\sigma$};
    \draw[->, thick] (0, -1.7) -- (0, 1.9) node[above, font=\footnotesize] {$j\omega$};
    \draw[dashed, thick, gray] (-2, -1.7) -- (-2, 1.7);
    \draw[dashed, thick, gray] (-1, -1.7) -- (-1, 1.7);

    \node[red, font=\normalsize] at (-1, 0) {$\times$};
    \node[red, font=\normalsize] at (-2, 0) {$\times$};
    \node[below, font=\scriptsize] at (-1, -0.12) {$-1$};
    \node[below, font=\scriptsize] at (-2, -0.12) {$-2$};

    \node[utnblue!80!black, font=\scriptsize] at (-1.5, 1.35) {ROC};

    % --- Conector L^{-1} ---
    \draw[->, gray, thick] (1.9, 0) -- (2.9, 0)
        node[midway, above, font=\scriptsize, gray] {$\mathcal{L}^{-1}$};

    % --- x(t): bilateral, negativa en ambos lados, salto en t = 0 ---
    \begin{scope}[shift={(4.8, 0)}]
        \draw[->, thick] (-1.7, 0) -- (2.6, 0) node[right, font=\footnotesize] {$t$};
        \draw[->, thick] (0, -1.8) -- (0, 0.6) node[right, font=\footnotesize] {$x(t)$};

        % Rama anticausal: -2e^{-t} sobre t < 0, factor 0.25
        \draw[utnblue!90!black, very thick, smooth, samples=60, domain=-1.2:0]
            plot (\x, {-0.5*exp(-\x)});
        \draw[utnblue!90!black, very thick, dashed, smooth, samples=20, domain=-1.27:-1.2]
            plot (\x, {-0.5*exp(-\x)});
        % Rama causal: -e^{-2t} sobre t > 0, mismo factor 0.25
        \draw[utnblue!90!black, very thick, smooth, samples=60, domain=0:2.4]
            plot (\x, {-0.25*exp(-2*\x)});
    \end{scope}

    \node[anchor=south, font=\small] at (2.0, 2.1) {(c)};
\end{scope}

\end{tikzpicture}
\end{document}
